\documentclass[11pt,a4paper]{article}
\usepackage[T1]{fontenc}
\usepackage[utf8]{inputenc}
\usepackage[margin=27mm,headheight=15pt]{geometry}
\usepackage{amsmath,amsthm,amssymb,mathtools}
\usepackage{newpxtext,newpxmath}
\usepackage{microtype}
\usepackage[table]{xcolor}
\usepackage{booktabs,array,longtable}
\usepackage{enumitem}
\usepackage{listings}
\usepackage{fancyhdr}
\usepackage{titlesec}
\usepackage{tocloft}
\usepackage{placeins}
\usepackage{tcolorbox}
\usepackage{hyperref}
\definecolor{ink}{HTML}{19334A}
\definecolor{accent}{HTML}{006B72}
\definecolor{soft}{HTML}{F0F6F7}
\hypersetup{colorlinks=true,linkcolor=ink,citecolor=accent,urlcolor=accent,
 pdftitle={Cycles refute infinite log-concavity for chromatic coefficients},
 pdfauthor={Research report prepared with ChatGPT}}
\titleformat{\section}{\large\bfseries\color{ink}}{\thesection}{0.7em}{}
\titleformat{\subsection}{\normalsize\bfseries\color{accent}}{\thesubsection}{0.7em}{}
\pagestyle{fancy}
\fancyhf{}
\fancyhead[L]{\small\color{ink}Chromatic coefficients and iterated log-concavity}
\fancyhead[R]{\small Research report}
\fancyfoot[C]{\small\thepage}
\renewcommand{\headrulewidth}{0.3pt}
\setlength{\parskip}{0.35em}
\setlength{\parindent}{1.1em}
\setcounter{tocdepth}{1}
\setlength{\cftbeforesecskip}{1pt}
\renewcommand{\cftsecfont}{\small\bfseries}
\renewcommand{\cftsecpagefont}{\small\bfseries}
\renewcommand{\cfttoctitlefont}{\large\bfseries\color{ink}}
\newtheorem{theorem}{Theorem}[section]
\newtheorem{proposition}[theorem]{Proposition}
\newtheorem{lemma}[theorem]{Lemma}
\newtheorem{corollary}[theorem]{Corollary}
\theoremstyle{definition}
\newtheorem{definition}[theorem]{Definition}
\newtheorem{remark}[theorem]{Remark}
\newcommand{\Lop}{\mathcal L}
\newcommand{\Z}{\mathbb Z}
\newcommand{\N}{\mathbb N}
\newcommand{\R}{\mathbb R}
\newcommand{\aco}{a^{(n)}}
\newcommand{\bin}[2]{\binom{#1}{#2}}
\lstset{language=Python,basicstyle=\ttfamily\footnotesize,
 keywordstyle=\color{ink}\bfseries,commentstyle=\color{accent},
 stringstyle=\color{accent},showstringspaces=false,breaklines=true,
 frame=single,rulecolor=\color{ink!25},columns=fullflexible,
 keepspaces=true,xleftmargin=0.7em,xrightmargin=0.7em}
\allowdisplaybreaks[2]
\emergencystretch=2em
\widowpenalty=10000
\clubpenalty=10000
\displaywidowpenalty=10000

\title{\vspace{-1.2em}\color{ink}\bfseries Cycles refute infinite log-concavity\\
for chromatic coefficients\\[0.55em]
\large Exact thresholds, finite certificates, and an exponential boundary limit}
\author{Research report prepared for Vladimir Reshetnikov\\
\small Prepared with ChatGPT; proofs and reproducible exact-arithmetic artifacts}
\date{September 19, 2026}

\begin{document}
\maketitle
\thispagestyle{empty}
\begin{abstract}
Amdeberhan and Moll's Conjecture~21 asserts that the absolute coefficients of
 every chromatic polynomial form an infinitely log-concave sequence. We give
 an elementary disproof using cycle graphs, and classify the entire cycle
 family. With zero padding and
 $\Lop(a)_k=a_k^2-a_{k-1}a_{k+1}$, the absolute coefficient sequence of $C_n$
 is infinitely log-concave exactly when $3\le n\le11$. Its first negative
 iterate occurs at depth $5$ for $n=12$, depth $4$ for $13\le n\le16$, and
 depth $3$ for all $n\ge17$. Every cycle sequence is nevertheless
 $2$-log-concave. The infinite negative family follows from a factored
 degree-$16$ polynomial; the positive cases are certified by an invariant
 $3$-factor log-concavity region. A normalized limit with coefficients of
 $e^x-1$ explains the failure and gives a perturbation-stable obstruction.
 We also derive generating functions and a constant-coefficient recurrence
 for the integer witness sequence. All finite checks use exact integers or
 rational numbers. The report does not claim priority over all existing
 literature, nor minimality among graphs other than cycles.
\end{abstract}

\begin{tcolorbox}[colback=soft,colframe=accent,boxrule=0.5pt]
\textbf{A compact disproof.} For $C_{17}$, the third iterate has
\[
  (\Lop^3\aco)_2
  =22\,491\,136^2-65\,536\cdot8\,150\,077\,440
  =-28\,272\,276\,537\,344<0.
\]
The complete local calculation appears in Section~\ref{sec:witness}; no
 floating-point calculation or exhaustive graph search is required.
\end{tcolorbox}

{\setlength{\parskip}{0pt}\small\tableofcontents}
\clearpage

\section{The problem, conventions, and main result}\label{sec:introduction}

\subsection{A precise published conjecture}
The target is Conjecture~21 in \emph{Infinite log-convexity} by
Amdeberhan and Moll~\cite{AM}: the absolute values of the coefficients of
any chromatic polynomial should be infinitely log-concave. The statement
appears on printed page~9 of the publisher's PDF; bibliographic and
manuscript-date details are recorded in the references.

This is a stronger assertion than ordinary log-concavity of chromatic
coefficients, established by Huh~\cite{Huh}. A different sufficient
condition comes from Br\"and\'en~\cite{Branden}: a polynomial with
nonnegative coefficients and only real nonpositive zeros has an infinitely
log-concave coefficient sequence. Neither theorem implies the proposed
statement for all chromatic polynomials.

\paragraph{Status and attribution.}
The primary paper inspected on September~19, 2026 still formulates the
assertion as a conjecture. A targeted search did not locate a prior cycle
counterexample, but it was not an exhaustive bibliographic search. Thus
this report establishes a disproof of the stated formulation, without
claiming to certify that it remained unresolved everywhere until this
report. The invariant-cone method used below is established machinery,
attributed to McNamara and Sagan~\cite{MS}; its proof is included.

\subsection{The nonlinear operator}
\begin{definition}
For a finite sequence $a=(a_0,\ldots,a_N)$, set $a_k=0$ when $k<0$ or
$k>N$, and define
\begin{equation}\label{eq:L}
  (\Lop a)_k=a_k^2-a_{k-1}a_{k+1}.
\end{equation}
A nonnegative sequence is \emph{$r$-log-concave} when $\Lop^j a$ is
entrywise nonnegative for every $1\le j\le r$. It is
\emph{infinitely log-concave} when this holds for every $j\ge1$.
\end{definition}
Zero extension is the convention explicitly used in~\cite{MS}. It keeps
both endpoints: the first and last nonzero entries are squared, rather
than discarded, at each iteration. In particular, we do not take absolute
values again after applying $\Lop$.

For a sequence under consideration, define its \emph{first failure depth}
by
\begin{equation}\label{eq:depth}
  h(a)=\min\{j\ge1:\text{$\Lop^j a$ has a negative entry}\},
\end{equation}
with $h(a)=\infty$ if this set is empty. Write $h(n)=h(\aco)$ for the
absolute coefficient sequence of the simple cycle $C_n$, indexed by
increasing degree.

\begin{theorem}[Complete cycle classification]\label{thm:main}
For every integer $n\ge3$, the absolute coefficients of the chromatic
polynomial of $C_n$ form the sequence
\begin{equation}\label{eq:row}
  a^{(n)}_0=0,\qquad a^{(n)}_1=n-1,\qquad
  a^{(n)}_k=\bin nk\quad(2\le k\le n).
\end{equation}
Every such sequence is $2$-log-concave, and
\begin{equation}\label{eq:classification}
 h(n)=
 \begin{cases}
   \infty,&3\le n\le11,\\
   5,&n=12,\\
   4,&13\le n\le16,\\
   3,&n\ge17.
 \end{cases}
\end{equation}
In every finite case, the coefficient of degree $2$ is negative at the
first failure depth. In particular, $C_{12}$ is the smallest cycle
counterexample to the infinite-log-concavity conjecture.
\end{theorem}

The qualification ``cycle'' in the last sentence matters. No exhaustive
search through all graphs on at most eleven vertices is asserted here.
The proof combines a symbolic argument valid for every $n\ge17$ with a
small collection of exact certificates for $3\le n\le16$.

\section{The cycle polynomial and its coefficient triangle}\label{sec:cycle}

\begin{proposition}\label{prop:chromatic}
For $n\ge3$, the chromatic polynomial and its absolute-coefficient
polynomial are respectively
\begin{align}
 P_{C_n}(q)&=(q-1)^n+(-1)^n(q-1),\label{eq:chromatic}\\
 A_n(x)&=\sum_{k=0}^n a^{(n)}_kx^k
       =(1+x)^n-(1+x).\label{eq:An}
\end{align}
\end{proposition}
\begin{proof}
The triangle has $q(q-1)(q-2)$ proper colorings, which agrees with
\eqref{eq:chromatic}. For $n\ge4$, delete one edge of $C_n$. The remaining
path has $q(q-1)^{n-1}$ proper colorings. The colorings in which the two
endpoints receive the same color are precisely those of the contracted
cycle $C_{n-1}$. Therefore
\[
  P_{C_n}(q)=q(q-1)^{n-1}-P_{C_{n-1}}(q).
\]
Substituting the induction hypothesis proves \eqref{eq:chromatic}.
The constant coefficient cancels, the linear coefficient has magnitude
$n-1$, and all coefficients of degree at least two are alternating
binomial coefficients. Equivalently,
$(-1)^nP_{C_n}(-x)=(1+x)^n-(1+x)$ has nonnegative coefficients, proving
\eqref{eq:An} and \eqref{eq:row}.
\end{proof}

Thus the sequences are almost rows of Pascal's triangle: their constant
coefficient is replaced by zero and their linear coefficient is reduced
by one. An arbitrarily small relative change in the latter coefficient
does not make the nonlinear iteration harmless, because the change at
the constant coefficient is not small.

The row polynomials satisfy two useful recurrences:
\begin{align}
 A_{n+1}(x)&=(1+x)A_n(x)+x(1+x),\label{eq:rowrec}\\
 A_n(x)&=A_{n-1}(x)+x(1+x)^{n-1}\qquad(n\ge4).\label{eq:delconrec}
\end{align}
The second is also a convenient independent way to generate the
coefficient triangle. At the coefficient level, \eqref{eq:rowrec} is
Pascal's recurrence with a correction of $1$ in degrees $1$ and $2$:
\[
 a^{(n+1)}_k=a^{(n)}_k+a^{(n)}_{k-1}
              +\mathbf 1_{k=1}+\mathbf 1_{k=2}.
\]
The bivariate ordinary generating function is rational:
\begin{equation}\label{eq:bivariate}
 \sum_{n\ge3} A_n(x)t^n
 =\frac{(1+x)^3t^3}{1-(1+x)t}
  -\frac{(1+x)t^3}{1-t}.
\end{equation}
These identities also specify the integer triangle in the bundled data
without relying on an external sequence identifier.

\paragraph{Indexing is not the obstruction.}
For a shift $T_s$ on zero-extended sequences, $(T_sa)_k=a_{k-s}$,
we have $\Lop T_s=T_s\Lop$. The operator also commutes with reversal of a
finite coefficient list. Consequently the counterexample persists when
coefficients are listed in decreasing degree or when the initial zero
coefficient is omitted and the remaining entries are reindexed. Deleting
endpoints at \emph{every} iteration, by contrast, would define a different
operator.

\section{A short exact counterexample: the seventeen-cycle}\label{sec:witness}

Only six input coefficients are needed to compute $(\Lop^3a)_2$.
For $C_{17}$ they are
\[
 (a_0,a_1,a_2,a_3,a_4,a_5)=(0,16,136,680,2380,6188).
\]
Let $b=\Lop a$, $c=\Lop^2a$, and $d=\Lop^3a$. Table~\ref{tab:triangle}
contains every nontrivial intermediate quantity needed for $d_2$.
Blank positions are unnecessary, not assumed to be zero.

\begin{table}[htbp]
\centering
\small
\begin{tabular}{@{}crrrrrr@{}}
\toprule
Iterate & $k=0$ & $k=1$ & $k=2$ & $k=3$ & $k=4$ & $k=5$\\
\midrule
$a$ & 0 & 16 & 136 & 680 & 2380 & 6188\\
$b$ & 0 & 256 & 7616 & 138720 & 1456560 & \\
$c$ & 0 & 65536 & 22491136 & 8150077440 & & \\
\bottomrule
\end{tabular}
\caption{An exact dependency triangle for the counterexample $C_{17}$.}
\label{tab:triangle}
\end{table}

For example,
\begin{align*}
 b_2&=136^2-16\cdot680=7616,\\
 b_3&=680^2-136\cdot2380=138720,\\
 b_4&=2380^2-680\cdot6188=1456560,\\
 c_2&=7616^2-256\cdot138720=22491136,\\
 c_3&=138720^2-7616\cdot1456560=8150077440.
\end{align*}
Since $a_0=b_0=c_0=0$, the first positive entry is repeatedly squared,
so $c_1=16^4=65536$. Finally,
\begin{align}
 d_2
  &=22491136^2-65536\cdot8150077440\notag\\
  &=505851198570496-534123475107840\notag\\
  &=-28272276537344<0.\label{eq:C17}
\end{align}
This exact calculation by itself disproves the universal conjecture. The
remaining sections explain why the example works, prove an infinite
family, and identify the smaller cycles that fail only at later depths.

\section{Symbolic iterates and the threshold at seventeen}\label{sec:symbolic}

Throughout this section, $a=\aco$, $b=\Lop a$, $c=\Lop^2a$, and
$d=\Lop^3a$. We compare them with a full Pascal row
$p_k=\bin nk$ and its successive iterates
$B=\Lop p$, $C=\Lop^2p$, and $D=\Lop^3p$.

\subsection{The first two Pascal iterates}
For $0\le k\le n$, direct cancellation of binomial ratios gives
\begin{align}
 B_k&=p_k^2\frac{n+1}{(k+1)(n-k+1)},\label{eq:B}\\
 C_k&=B_k^2\frac{2(n+2)}{(k+2)(n-k+2)}.\label{eq:C}
\end{align}
For interior indices, the corresponding neighboring-product ratios are
\begin{align}
 \frac{B_{k-1}B_{k+1}}{B_k^2}
   &=\frac{k(n-k)}{(k+2)(n-k+2)},\label{eq:B-ratio}\\
 \frac{C_{k-1}C_{k+1}}{C_k^2}
   &=\frac{k^2(n-k)^2}
    {(k+1)(k+3)(n-k+1)(n-k+3)}.\label{eq:C-ratio}
\end{align}
Both ratios are strictly less than one. At either endpoint the
neighboring product is zero, and the same positivity conclusion holds.
In particular $B_k,C_k,D_k>0$ for $0\le k\le n$. These elementary
identities suffice for the tail comparisons below; no general
real-rootedness theorem is needed in the proof of the cycle classification.

\subsection{The exceptional coefficients}
Because $a_0=0$, we have $b_0=c_0=d_0=0$. The first iterate begins with
\begin{align}
 b_1&=(n-1)^2,\label{eq:b1}\\
 b_2&=\frac{n(n-1)^2(n+4)}{12},\label{eq:b2}\\
 b_k&=B_k\qquad(3\le k\le n).\label{eq:b-tail}
\end{align}
For later calculations the two specific tail values are
\begin{align}
 b_3&=\frac{n^2(n-1)^2(n-2)(n+1)}{144},\label{eq:b3}\\
 b_4&=\frac{n^2(n-1)^2(n-2)^2(n-3)(n+1)}{2880}.
 \label{eq:b4}
\end{align}
The formula for $b_4$ is also valid at $n=3$, where it is zero by padding.
All entries on the positive support of $b$ are positive.

The second iterate has
\begin{align}
 c_1&=(n-1)^4,\label{eq:c1}\\
 c_2&=\frac{n^2(n-1)^4(n+2)}{16},\label{eq:c2}\\
 c_3&=\frac{n^3(n-2)^2(n-1)^4(n+1)(n^2+n+18)}{51840},
       \label{eq:c3}\\
 c_k&=C_k\qquad(4\le k\le n).\label{eq:c-tail}
\end{align}
For instance, $c_2=b_2^2-b_1b_3$ reduces to
\[
 \frac{n^2(n-1)^4}{144}
       \bigl((n+4)^2-(n-2)(n+1)\bigr)
 =\frac{n^2(n-1)^4(n+2)}{16}.
\]
Substitution of \eqref{eq:b2}--\eqref{eq:b4} gives \eqref{eq:c3}.
The exact algebra is also audited by \texttt{derive\_symbolic.py}.

\begin{proposition}\label{prop:2LC}
For every $n\ge3$, the cycle coefficient sequence is $2$-log-concave.
In fact, both its first and second iterates are strictly positive on
indices $1,\ldots,n$.
\end{proposition}
\begin{proof}
Every factor in \eqref{eq:b1}--\eqref{eq:b-tail} and
\eqref{eq:c1}--\eqref{eq:c-tail} is positive on the indicated support.
The remaining entries are zero. The positivity of the Pascal terms was
proved using \eqref{eq:B} and \eqref{eq:C}.
\end{proof}

\subsection{The decisive quartic}
Define
\begin{equation}\label{eq:Q}
 Q(n)=2n^4-12n^3-327n^2-412n+36.
\end{equation}
Using $d_2=c_2^2-c_1c_3$, we obtain
\begin{equation}\label{eq:d2}
 \boxed{\displaystyle
 d_2=-\frac{n^3(n-1)^8(n+4)}{103680}\,Q(n).}
\end{equation}
An explicit intermediate form, useful for checking the factorization, is
\[
 d_2=\frac{n^3(n-1)^8}{103680}
 \left(405n(n+2)^2
  -2(n-2)^2(n+1)(n^2+n+18)\right).
\]
The polynomial in parentheses is exactly $-(n+4)Q(n)$.

\begin{lemma}\label{lem:Qsign}
For integral $n\ge3$, $Q(n)<0$ when $n\le16$, and $Q(n)>0$ when $n\ge17$.
\end{lemma}
\begin{proof}
On $3\le n\le16$, the increasing quadratic
$2n^2-12n-327$ is at most its value $-7$ at $n=16$. Hence
\[
 Q(n)=n^2(2n^2-12n-327)-412n+36<0.
\]
For $n=17+m$, $m\ge0$, expansion gives
\begin{equation}\label{eq:Qshift}
 Q(17+m)=2m^4+124m^3+2529m^2+17370m+6615>0.
\end{equation}
\end{proof}

\begin{corollary}\label{cor:infinitefamily}
For every $n\ge17$, $\aco$ fails at its third iterate. In particular,
the conjecture has infinitely many connected, simple counterexamples.
\end{corollary}
\begin{proof}
The prefactor outside $Q(n)$ in \eqref{eq:d2} is negative and nonzero.
Lemma~\ref{lem:Qsign} therefore gives $d_2<0$. Proposition~\ref{prop:2LC}
shows that no earlier failure is possible.
\end{proof}

\subsection{No other third-iterate obstruction for smaller cycles}
To prove exact failure depths for $12\le n\le16$, it remains to know
that their \emph{entire} third iterates are nonnegative, not just $d_2$.
Define
\[
 S(n)=4n^5-27n^4+186n^3+145n^2+684n-3564.
\]
Direct calculation gives
\begin{equation}\label{eq:d3}
 d_3=\frac{n^6(n-2)^3(n-1)^8(n+1)^2}{10749542400}\,S(n).
\end{equation}
Every $n\ge3$ has $S(n)>0$, since
\begin{equation}\label{eq:Sshift}
 S(3+m)=4m^5+33m^4+222m^3+1441m^2+5280m+3600.
\end{equation}
The next coordinate admits a useful comparison rather than a large
factorization. From the first-iterate formulas,
\[
  b_2=B_2+p_3,
  \qquad c_3=C_3-p_3B_4\le C_3.
\]
Since $c_k=C_k$ for $k\ge4$, it follows that
\begin{equation}\label{eq:d4compare}
 d_4=C_4^2-c_3C_5\ge C_4^2-C_3C_5=D_4>0
 \quad(4\le n),
\end{equation}
where $C_5=0$ when $n=4$. For $k\ge5$ on the support, $d_k=D_k>0$.
Finally $d_1=(n-1)^8>0$.

\begin{proposition}\label{prop:3LC}
The cycle coefficient sequence is $3$-log-concave exactly when
$3\le n\le16$. In that range its third iterate is strictly positive
on the nonzero support.
\end{proposition}
\begin{proof}
The formulas above establish positivity at every index other than $2$;
Lemma~\ref{lem:Qsign} and \eqref{eq:d2} settle index $2$. For $n=3$,
indices beyond $3$ are simply absent. Proposition~\ref{prop:2LC} supplies
the earlier two iterates.
\end{proof}

\section{Infinite positivity for the nine small cycles}\label{sec:cone}

A long list of positive iterates is not a proof of infinite
log-concavity. Instead we use a region that $\Lop$ maps into itself.
The following sufficient criterion is the McNamara--Sagan
invariant-cone argument~\cite[Lemma~2.1]{MS}; using the rational constant
$3$ makes the certificates especially simple.

\begin{lemma}[Invariant factor-log-concavity region]\label{lem:cone}
Let $r\ge(3+\sqrt5)/2$ and let $a$ be nonnegative with
\begin{equation}\label{eq:rcone}
 a_k^2\ge r\,a_{k-1}a_{k+1}\qquad\text{for all }k.
\end{equation}
Then $b=\Lop a$ is nonnegative and satisfies the same inequalities.
Consequently $a$ is infinitely log-concave. In particular, $r=3$ is
sufficient.
\end{lemma}
\begin{proof}
The assumption gives, coordinatewise,
\[
 (1-1/r)a_k^2\le b_k\le a_k^2,
 \qquad b_k\ge0.
\]
Therefore
\begin{align*}
 b_k^2
 &\ge (1-1/r)^2a_k^4\\
 &\ge (r-1)^2a_{k-1}^2a_{k+1}^2\\
 &\ge r\,b_{k-1}b_{k+1}.
\end{align*}
The last inequality uses $(r-1)^2\ge r$ and
$0\le b_{k\pm1}\le a_{k\pm1}^2$. All steps remain valid when a coordinate
is zero; no division by an entry of $a$ is used. Iteration proves the claim.
\end{proof}

For a nonnegative vector $u$ with at least one positive neighboring
product, let
\[
 R(u)=\min_{k:\,u_{k-1}u_{k+1}>0}
        \frac{u_k^2}{u_{k-1}u_{k+1}}.
\]
If $R(u)\ge3$, indices with zero neighboring product satisfy the cone
condition automatically. Table~\ref{tab:cone} gives the exact minimum
for a suitable iterate of each small cycle.

\begin{table}[htbp]
\centering
\begin{tabular}{@{}rrrl@{}}
\toprule
$n$ & Iterate $j$ & Minimizing index $k$ & $R(\Lop^j\aco)$\\
\midrule
3 & 0 & 2 & $9/2$\\
4 & 1 & 3 & $25/6$\\
5 & 2 & 3 & $256/49$\\
6 & 2 & 3 & $100/27$\\
7 & 2 & 3 & $6845/2187$\\
8 & 3 & 3 & $30926/5863$\\
9 & 3 & 3 & $8736525/1666093$\\
10 & 3 & 3 & $2991871581875/519958761651$\\
11 & 3 & 2 & $32433025/8287488$\\
\bottomrule
\end{tabular}
\caption{Finite certificates of infinite log-concavity. Every displayed
ratio is greater than $3$. All ratios and minima are exact.}
\label{tab:cone}
\end{table}

The entries are obtained from \eqref{eq:row} by applying \eqref{eq:L}
at most three times and comparing fractions by cross multiplication.
The complete vectors are provided in \texttt{data/certificates.json};
Appendix~\ref{app:code} gives a short, independent-of-SymPy checker of
all these minima. Thus each row represents finitely many explicit integer
inequalities, not an extrapolation from a numerical plot.

\begin{proposition}\label{prop:small}
The sequences $\aco$ are infinitely log-concave for $3\le n\le11$.
\end{proposition}
\begin{proof}
All iterates preceding those in Table~\ref{tab:cone} are nonnegative by
Propositions~\ref{prop:2LC} and \ref{prop:3LC}, or by the initial row.
The tabulated iterate is nonnegative and satisfies the $r=3$ condition.
Lemma~\ref{lem:cone} certifies every subsequent iterate.
\end{proof}

It is worth distinguishing a certificate from a necessary condition:
this proof does not assert that every infinitely log-concave sequence
eventually enters this cone. Only the nine displayed cases require that
property here, and for them it is checked exactly.

\section{The delayed failures at twelve through sixteen}\label{sec:late}

For $12\le n\le16$, Proposition~\ref{prop:3LC} shows that depth three is
still positive. The next calculations determine the exact first
failure. In Table~\ref{tab:failures}, the middle column is the first
failure depth, and the final column is the offending entry of that
iterate.

\begin{table}[htbp]
\centering
\small
\begin{tabular}{@{}rrr@{}}
\toprule
$n$ & $h(n)$ & $(\Lop^{h(n)}\aco)_2$\\
\midrule
12 & 5 & $-249621701601023742801969101519201265582387397744$\\
13 & 4 & $-3618341131654935620812800$\\
14 & 4 & $-199158562975246657489530096$\\
15 & 4 & $-2734560032125157358883149375$\\
16 & 4 & $-25982668618402950000000000000$\\
17 & 3 & $-28272276537344$\\
\bottomrule
\end{tabular}
\caption{Exact negative witnesses. The $n=17$ row is repeated for
comparison with the delayed failures.}
\label{tab:failures}
\end{table}

For $n=13,14,15,16$, one can reproduce the values without constructing a
full fourth-iterate vector: use
\begin{equation}\label{eq:e2}
  (\Lop^4a)_2=d_2^2-d_1d_3
\end{equation}
with $d_1=(n-1)^8$ and the formulas \eqref{eq:d2}, \eqref{eq:d3}.
Their earlier iterates are already known to be positive, so these
negative values settle their exact depths.

For $n=12$, all fourth-iterate entries are nonnegative. The complete
fourth-iterate vector, with every integer explicitly displayed, appears
in Appendix~\ref{app:C12}. Its first three positive entries are
\begin{align*}
 e_1&=45949729863572161,\\
 e_2&=1505318939996586804452286,\\
 e_3&=54746933663864342589859620161140.
\end{align*}
Consequently
\[
 (\Lop^5a^{(12)})_2=e_2^2-e_1e_3
  =-249621701601023742801969101519201265582387397744<0.
\]
This also explains why checking only three or four iterations would miss
some counterexamples.

\begin{proof}[Completion of Theorem~\ref{thm:main}]
Proposition~\ref{prop:chromatic} gives the coefficient row;
Proposition~\ref{prop:2LC} gives universal $2$-log-concavity;
Corollary~\ref{cor:infinitefamily} gives $h(n)=3$ for $n\ge17$;
Proposition~\ref{prop:small} gives the nine infinite cases; and the exact
calculations in this section give all remaining depths. These cases
exhaust $n\ge3$.
\end{proof}

\section{Why the failure is structural: an exponential limit}\label{sec:limit}

\subsection{Locality and geometric normalization}
For $t>0$, define the geometric rescaling $(G_ta)_k=t^ka_k$. Directly
from \eqref{eq:L},
\begin{equation}\label{eq:scaling}
 \Lop^j(G_ta)_k=t^{2^j k}(\Lop^ja)_k.
\end{equation}
This follows by induction, since each iteration squares the scale.
Moreover, $(\Lop^ja)_k$ is a polynomial in the entries
$a_{k-j},\ldots,a_{k+j}$, with negative-index entries fixed at zero.
Each iteration enlarges the dependency interval by one on each side.

Taking $t=1/n$ gives a normalized row $u^{(n)}_k=a^{(n)}_k/n^k$. For
fixed $k$,
\[
 u^{(n)}_0=0,\qquad u^{(n)}_1\longrightarrow1,\qquad
 u^{(n)}_k\longrightarrow\frac1{k!}\quad(k\ge2).
\]
Thus the coefficientwise limit is the sequence $u$ of $e^x-1$:
\begin{equation}\label{eq:expcoeff}
 u=(0,1,1/2!,1/3!,1/4!,1/5!,\ldots).
\end{equation}

The necessary local iterates of this limiting sequence are particularly
simple:
\begin{align*}
 u&=(0,1,1/2,1/6,1/24,1/120,\ldots),\\
 \Lop u&=(0,1,1/12,1/144,1/2880,\ldots),\\
 \Lop^2u&=(0,1,0,1/51840,\ldots).
\end{align*}
The zero between two positive entries forces
\begin{equation}\label{eq:limitnegative}
   (\Lop^3u)_2=-\frac1{51840}.
\end{equation}
Only a finite prefix is involved, so passage to the limit follows from
continuity of a polynomial; there is no interchange of an infinite number
of nonlinear iterations with a limit.

\begin{theorem}[A local, perturbation-stable obstruction]\label{thm:local}
Let $v^{(N)}$ be nonnegative sequences indexed by the nonnegative integers,
zero-extended to negative indices. Suppose $s_N>0$ and
\[
 v^{(N)}_0\longrightarrow0,\qquad
 \frac{v^{(N)}_k}{s_N^k}\longrightarrow\frac1{k!}
       \quad(1\le k\le5).
\]
Then
\begin{equation}\label{eq:local-limit}
 s_N^{-16}(\Lop^3v^{(N)})_2\longrightarrow-\frac1{51840}.
\end{equation}
In particular, all sufficiently late members of the family fail
$3$-log-concavity.
\end{theorem}
\begin{proof}
Apply \eqref{eq:scaling} with $t=1/s_N$ and $j=3,k=2$.
The required entry depends only on indices $-1$ through $5$.
Those entries converge to the corresponding ones in \eqref{eq:expcoeff},
with the negative index fixed at zero. Polynomial continuity and
\eqref{eq:limitnegative} prove the limit. Its strict negativity gives an
eventually negative entry.
\end{proof}

For example, for any fixed real $\lambda$, the rows
\[
 (0,n-\lambda,\bin n2,\ldots,\bin nn)
\]
are nonnegative for all sufficiently large $n$ and fail $3$-log-concavity
eventually. More generally, the same conclusion survives arbitrary
perturbations $o(n^k)$ in coordinates $1\le k\le5$, provided the zeroth
coordinate tends to zero and the resulting rows remain nonnegative.
This is a statement about local coefficient data, not about whether every
such perturbation is realized by a chromatic polynomial.

\subsection{A quantitative asymptotic expansion}
Equation~\eqref{eq:d2} sharpens the limiting argument to an exact
polynomial. Expanding its highest powers gives
\begin{equation}\label{eq:asymptotic}
 (\Lop^3\aco)_2
 =-\frac{n^{16}}{51840}
   +\frac{n^{15}}{5184}
   +\frac{287n^{14}}{103680}
   +O(n^{13}).
\end{equation}
The remainder is simply the lower-degree part of a polynomial, not an
uncontrolled approximation. The asymptotic limit explains eventual
failure, whereas the positive-shift quartic \eqref{eq:Qshift} proves the
exact threshold $17$.

The mechanism is a boundary effect. The full Pascal row, after the same
normalization, tends to the coefficients of $e^x$, whose constant term
is $1$. Our rows instead tend to $e^x-1$. The changed boundary produces a
zero in the second iterate and a strict negative entry in the third.
It would therefore be unsafe to transfer infinite log-concavity from
Pascal rows merely because most coefficients are unchanged.

\section{An integer witness sequence, its recurrence and generating function}
\label{sec:witnesssequence}

The negative obstruction itself defines a natural positive integer
sequence. Set
\begin{equation}\label{eq:wm}
 w_m=-(\Lop^3a^{(17+m)})_2\qquad(m\ge0).
\end{equation}
From \eqref{eq:d2} and \eqref{eq:Qshift},
\begin{equation}\label{eq:wmfactor}
 \begin{split}
 w_m={}&\frac{(m+17)^3(m+16)^8(m+21)}{103680}\\
 &\mathrel{}\times
 (2m^4+124m^3+2529m^2+17370m+6615).
 \end{split}
\end{equation}
Integrality follows directly from the integer sequence recurrence defining
$\Lop$; it need not be inferred from the rational prefactor.

The first three values are
\[
 28272276537344,\quad229969795557447,\quad880220808750663.
\]
Equation~\eqref{eq:wmfactor} expresses $w_m$ as a degree-$16$ polynomial
in $m$, with strictly positive coefficients and leading coefficient
$1/51840$. In particular it is strictly increasing for $m\ge0$.

Let $\Delta f(m)=f(m+1)-f(m)$ and let $\delta_j=\Delta^jw_0$.
Then Newton's finite expansion gives
\begin{equation}\label{eq:newton}
   w_m=\sum_{j=0}^{16}\delta_j\bin mj.
\end{equation}
Every $\delta_j$ for $0\le j\le16$ is a positive integer. To see
positivity, express the polynomial in powers $m^r$: the coefficients
are positive, and $\Delta^jm^r|_{m=0}=j!\,{r\brace j}\ge0$.
For $j=0$ the positive constant term suffices; for $j\ge1$ the $m^j$
coefficient already contributes positively.

Consequently the ordinary generating function is
\begin{equation}\label{eq:wgf}
 W(t)=\sum_{m\ge0}w_mt^m
     =\sum_{j=0}^{16}\frac{\delta_jt^j}{(1-t)^{j+1}}
     =\frac{P(t)}{(1-t)^{17}},
\end{equation}
where an exact specification of the degree-at-most-$16$ numerator is
\begin{equation}\label{eq:Pnumerator}
 P(t)=\sum_{j=0}^{16}\delta_jt^j(1-t)^{16-j},\qquad
 \delta_j=\sum_{i=0}^j(-1)^{j-i}\bin ji w_i.
\end{equation}
These are finite closed formulas, rather than auxiliary recurrences
needed to define the coefficients. The integer lists $\delta_j$ and
$[t^j]P(t)$ are included in
\path{data/witness_generating_function.json}.

The denominator has a pole of exact order $17$ at $t=1$, since
$P(1)=\delta_{16}=16!/51840>0$. Equivalently, the sequence satisfies
\begin{equation}\label{eq:wrec}
 \sum_{j=0}^{17}(-1)^{17-j}\bin{17}{j}w_{m+j}=0
 \qquad(m\ge0).
\end{equation}
Thus the original graph-theoretic question yields an explicit integer
sequence with a rational generating function, a constant-coefficient
recurrence, and a complete polynomial asymptotic expansion. No OEIS
accession or novelty claim for this derived sequence is assumed.

\section{Verification, reproducibility, and proof boundaries}\label{sec:verification}

The proof has three logically different components. The explicit
$C_{17}$ arithmetic already disproves the conjecture. The symbolic
factorization and its sign proof treat every $n\ge17$. Finally, a finite
set of rational inequalities and integer identities settles the cases
$3\le n\le16$. Only the last component needs finite certificate
checking, and those certificates are fully specified in the report and
accompanying files.

\subsection{What was executed}
The standard-library program \texttt{verify.py} was run with full cycle
rows through $n=200$. It checks the first-failure classification, all nine
positive cone certificates, the formulas for the third iterate, and an
independent deletion--contraction generation of the coefficient rows.
It also computes the necessary six-entry prefix at
$n=1000,10000,10^6,10^{12}$ and compares the exact third-iterate result
with \eqref{eq:d2}. These larger checks are corroboration, not the proof
of the universal assertion.

A separate SymPy program, \texttt{derive\_symbolic.py}, was executed to
audit all displayed low-index factorizations, both positive-shift
identities, the Pascal neighbor ratios, the asymptotic coefficients, and
the rational generating-function data. The recorded environment was
Python~3.13.5 and SymPy~1.14.0. The principal checker needs only the Python
standard library; SymPy is optional. No Lean formalization or independent
human referee verification is claimed.

\subsection{Exact arithmetic and computational cost}
A whole-row iteration uses $O(n)$ integer operations and $O(n)$ integer
storage if only the current row is retained. If entries initially have
at most $B$ bits, the largest bit length after $j$ iterations is
$O(2^j(B+1))$, because a new value is a difference of products or squares.
For a fixed depth, full-row tests are consequently practical at modest
$n$, but integer sizes grow rapidly with depth.

The local witness requires only six initial entries, three shrinking
prefix updates, and a constant number of integer operations after those
entries are formed. Their bit lengths are $O(\log n)$, and the witness
has $O(\log n)$ bits as well, since its degree is fixed. This explains
why the $n=10^{12}$ prefix test is easy without ever constructing a row
of length $10^{12}$.

Floating-point ratio iteration is unnecessary and can be misleading
near cancellation. The programs therefore neither convert large
coefficients to machine floats nor decide a sign from a decimal
approximation. They also never apply absolute values to an iterate:
once a negative entry appears, the infinite-log-concavity property has
already failed.

\subsection{Rebuilding the package}
From the archive directory, run
\begin{lstlisting}[language=bash]
python verify.py
python minimal_verifier.py
python derive_symbolic.py   # optional: needs SymPy
pdflatex -interaction=nonstopmode -halt-on-error article.tex
pdflatex -interaction=nonstopmode -halt-on-error article.tex
\end{lstlisting}
The \texttt{--write-data} option regenerates the JSON and CSV artifacts.
Sequence values in the JSON certificate are decimal strings, deliberately
avoiding precision loss in software that parses JSON numbers as binary
floating-point values. The CSV data are exact decimal integers.

The archive contains the article source and PDF, the three verification
programs, execution transcripts, a README, a build script, and the
coefficient and witness datasets. It contains no third-party font files
or copied research-paper PDFs.

\section{Consequences and remaining directions}\label{sec:consequences}

The conjecture fails within a particularly restrictive class: simple,
connected, $2$-regular graphs. The failure is therefore not a consequence
of loops, multiple edges, dense graphs, or complicated graph
constructions. On the other hand, it is genuinely a higher-iteration
phenomenon: every cycle passes the first two iterations.

Adding $s$ isolated vertices to a graph multiplies its chromatic
polynomial by $q^s$. It only shifts the coefficient sequence, and the
shift invariance of $\Lop$ preserves all failure depths. This supplies
further counterexamples, although the cycles already give a connected
infinite family. No analogous inheritance from an arbitrary subgraph is
claimed: the presence of a long cycle inside a larger graph does not by
itself transfer this coefficient calculation.

The positive cases illustrate the limits of the real-rooted sufficient
criterion. By \eqref{eq:An},
\[
 A_n(x)=(1+x)\bigl((1+x)^{n-1}-1\bigr).
\]
For $n\ge4$ this has nonreal roots obtained from nonreal $(n-1)$st roots
of unity by subtracting $1$. Nevertheless $C_4,\ldots,C_{11}$ are
infinitely log-concave by their cone certificates. Real-rootedness is
therefore sufficient but not necessary in this family.

Several questions are separate from what has been proved. The smallest
counterexample among \emph{all} simple graphs could be smaller than
$C_{12}$. Whether every chromatic coefficient sequence is $2$-log-concave
is not decided by the cycle calculations. One may also ask for a useful
structural description of those graphs whose coefficient sequences are
infinitely log-concave. These are directions for further investigation,
not claims of newly verified open-problem status.

The central methodological lesson is concrete: normalize a tractable
family, inspect a few boundary coefficients under the nonlinear
operator, and use exact arithmetic to expose a sign change. Here that
approach leads beyond a single counterexample to an exact classification
and a stable asymptotic obstruction.

\appendix
\section{The fourth iterate for the twelve-cycle}\label{app:C12}

For $a=a^{(12)}$, the vector $e=\Lop^4a$ has the entries in
Table~\ref{tab:C12}. All are nonnegative, and every entry of indices
$1$ through $12$ is positive. Together with
Proposition~\ref{prop:3LC}, this certifies that the negative fifth-iterate
entry in Section~\ref{sec:late} really is the first failure.

\begin{table}[htbp]
\centering
\small
\begin{tabular}{@{}rr@{}}
\toprule
$k$ & $e_k=(\Lop^4a^{(12)})_k$\\
\midrule
0 & 0\\
1 & 45949729863572161\\
2 & 1505318939996586804452286\\
3 & 54746933663864342589859620161140\\
4 & 15531438327857548927780352470984421505\\
5 & 11178866851263718584359046258900106285056\\
6 & 102599969549662636115648080319674359475968\\
7 & 11246650538385145317068550658400941510656\\
8 & 13507131422281444862413759705797857505\\
9 & 131200081206735639046019512080640\\
10 & 5555828662478863428996336\\
11 & 299028575209728\\
12 & 1\\
\bottomrule
\end{tabular}
\caption{A complete exact certificate that the fourth iterate for
$C_{12}$ has no negative entries.}
\label{tab:C12}
\end{table}

\FloatBarrier
\section{A short checker for the finite certificates}\label{app:code}

The following complete Python program computes the finite cases directly
from binomial coefficients. It requires only the standard library and
prints the exact minimizing ratios or negative entries. It uses no
precomputed iterate vectors. The larger \texttt{verify.py} adds
explicit expected values, data generation, redundant identity tests,
and checks that remain enabled under Python's optimized mode.

\lstinputlisting{minimal_verifier.py}

This finite program is not being asked to establish a statement by
checking infinitely many iterates. Its positive exit condition invokes
Lemma~\ref{lem:cone}; its negative exit condition provides an explicit
counterexample. That distinction is the reason the finite certificates
are logically sufficient.

\begin{thebibliography}{9}

\bibitem{AM}
T. Amdeberhan and V. H. Moll,
\emph{Infinite log-convexity},
Online Journal of Analytic Combinatorics, issue 17 (2022), article 6,
10 pages. Publisher PDF dated December 30, 2023; Conjecture 21 on p.~9.
\url{https://combinatorialpress.com/article/ojac/vol17/305.pdf}.
Accessed September 19, 2026.

\bibitem{MS}
P. R. W. McNamara and B. E. Sagan,
\emph{Infinite log-concavity: developments and conjectures},
arXiv:0808.1065v2 (2009), especially Section~1 and Lemma~2.1.
\url{https://arxiv.org/abs/0808.1065}.

\bibitem{Branden}
P. Br\"and\'en,
\emph{Iterated sequences and the geometry of zeros},
Journal f\"ur die reine und angewandte Mathematik \textbf{658} (2011),
115--131; arXiv:0909.1927.
\url{https://arxiv.org/abs/0909.1927}.

\bibitem{Huh}
J. Huh,
\emph{Milnor numbers of projective hypersurfaces and the chromatic
polynomial of graphs},
Journal of the American Mathematical Society \textbf{25} (2012),
907--927; arXiv:1008.4749.
\url{https://arxiv.org/abs/1008.4749}.

\end{thebibliography}
\end{document}
